\section{Safety Discussion}
\label{sec:safety}
% \vspace{-2mm}

Systems capable of self-improvement, such as the DGM, represent a step toward more autonomous AI development, aligning with long-standing goals in the field of making capable AI that can benefit humanity~\citep{schmidhuber1987evolutionary, clune2019ai, markoff2016machines, lehman2023machine}. However, this capability introduces unique safety considerations stemming from the system's ability to autonomously modify its own code. Modifications optimized solely for benchmark performance might inadvertently introduce vulnerabilities or behaviors misaligned with human intentions, even if they improve the target metric~\citep{bostrom2020ethical}. In particular, if evaluation benchmarks do not fully capture all desired agent properties (e.g., safety and robustness), the self-improvement loop could amplify misalignment over successive generations. Iterative self-modification could also lead to increasingly complex and uninterpretable internal logic, hindering human understanding, oversight, and control~\citep{sheth2025safety, anwar2024foundational, greenblatt2024alignment, ganguli2022red}.

Recognizing these challenges, the current implementation and experimental setup of the DGM incorporates several safeguards. All agent execution and self-modification processes are conducted within isolated sandboxed environments, limiting their ability to affect the host system, and thereby mitigating the risk of unintended actions. Each execution within the sandbox is subjected to a strict time limit, reducing the risk of resource exhaustion or unbounded behavior. The self-improvement process is currently confined to the well-defined domain of enhancing performance on specific coding benchmarks by modifying the agent's own Python codebase, thus limiting the scope of potential modifications. Additionally, we actively monitor agent performance and code changes, with the DGM archive providing a traceable lineage of modifications for review. At this stage, we have found no evidence of harmful or malicious behavior in the generated agents, and the self-modifications have been primarily focused on improving coding capabilities.

Conversely, a significant potential benefit of the self-improvement paradigm is that it could, in principle, be directed toward enhancing safety and interpretability themselves. We conduct a preliminary investigation into how the DGM can be deployed in AI safety settings to develop countermeasures for FM hallucination (\Cref{app:dgm-halluc}). Just as the DGM learns to improve its coding capabilities, it could potentially discover and integrate better internal safeguards or modify itself for greater transparency (e.g., incorporating principles akin to Constitutional AI~\citep{bai2022constitutional}), if such properties were included in its evaluation criteria~\citep{rosser2025agentbreeder}. This suggests a promising, albeit challenging, pathway in which self-improvement becomes a tool for building more trustworthy AI systems. Additional research could also explore weaving Constitutional AI in from the start, though the challenge would be incentivizing the system to retain these directives (an option worth exploring is to create an unmodifiable part of the system to be able to evaluate at halt the rest).

The DGM demonstrates the potential of self-improving AI while still operating within safe research boundaries due to the current limitations of frontier FMs and effective mitigations like sandboxing. \Cref{app:safety} presents additional discussion on broader safety uncertainties. We include this safety discussion proactively to raise awareness about the emerging prospect of self-improving AI systems and their associated safety implications, particularly as these systems inevitably become more capable~\citep{yudkowsky2008artificial, bostrom2002a,ecoffet2020open, bengio2024managing,clune2019ai}. Accordingly, we advocate for continued investigation into the safe and beneficial evolution of AI-Generating Algorithms~\citep{clune2019ai} and self-improving systems.

% \vspace{-2mm}
